% === Báo cáo mở rộng cho PoseAlert ===
\documentclass[12pt,a4paper]{article}
\usepackage[utf8]{vietnam}
\usepackage{amsmath,amssymb}
\usepackage{graphicx}
\usepackage{hyperref}
\usepackage{listings}
\usepackage{xcolor}
\usepackage{geometry}
\usepackage{longtable}
\usepackage{tabularx}
\usepackage{setspace}
\usepackage{caption}
\usepackage{float}
\usepackage{enumitem}
\geometry{a4paper, margin=1in}
\onehalfspacing

% Code listing style
\lstset{basicstyle=\ttfamily\footnotesize,breaklines=true}

\title{%
    \vspace{1.5cm}
    \textbf{\Huge BÁO CÁO DỰ ÁN}\\[0.3cm]
    \Large \textbf{Ứng dụng Web Cảnh Báo Tư Thế Ngồi Học — PoseAlert}\\[0.3cm]
    \large Ứng dụng nhận diện tư thế thời gian thực với TensorFlow.js và MoveNet
}
\author{Nhóm thực hiện: \begin{tabular}{l}Vương Lâm \\ [Tên thành viên khác — điền vào]\end{tabular}}
\date{\today}

\begin{document}

\maketitle
\thispagestyle{empty}
\newpage

\tableofcontents
\newpage

% -----------------------------
\section{Tóm tắt dự án}
\paragraph{}
PoseAlert là một ứng dụng web hướng tới mục tiêu giúp người dùng cải thiện tư thế ngồi khi học hoặc làm việc trước máy tính. Ứng dụng hoạt động hoàn toàn trên trình duyệt, sử dụng mô hình ước lượng tư thế (pose estimation) để trích xuất tọa độ các khớp người, sau đó áp dụng các luật hình học để phát hiện các trạng thái tư thế không phù hợp (ví dụ: cúi đầu, vẹo lưng, mắt quá gần). Bản báo cáo này mô tả toàn bộ quá trình thực hiện dự án: phân công công việc, thiết kế kỹ thuật, triển khai, kiểm thử, kết quả, điểm mạnh, khó khăn và kế hoạch phát triển tiếp theo.

\vspace{0.5cm}
\section{Danh sách thành viên và phân công}
\begin{longtable}{p{4.5cm} p{10cm}}
\textbf{Thành viên} & \textbf{Vai trò \& Nhiệm vụ chính} \\\hline
Vương Lâm & Quản lý dự án, Thiết kế giao diện (UI/UX), Triển khai frontend (HTML/CSS/JS), Tích hợp mô-đun nhận diện tư thế và smoothing logic. \\
Hoàng Trung Dũng & Kỹ sư backend (Firebase integration), Thiết kế hệ thống phòng học ảo, Quản lý lưu trữ, Chức năng chat và đồng bộ thời gian thực. \\
\end{longtable}

\section{Quá trình thực hiện (Timeline)}
\begin{itemize}[leftmargin=*]
    \item \textbf{Pha 1 (Nghiên cứu):} Tìm hiểu MoveNet, TensorFlow.js và lên ý tưởng tính năng.
    \item \textbf{Pha 2 (Thiết kế UI/UX):} Xây dựng bố cục 2-panel, module CSS, gamification.
    \item \textbf{Pha 3 (Tích hợp AI):} Xây dựng vòng lặp nhận diện và tinh chỉnh thuật toán.
    \item \textbf{Pha 4 (Tính năng phụ):} Thêm Pomodoro, Dashboard thống kê, hệ thống Nhiệm vụ (Quest).
    \item \textbf{Pha 5 (Realtime \& Firebase):} Tích hợp Chat, Phòng ảo, Gửi file qua Firebase.
    \item \textbf{Pha 6 (Kiểm thử):} Tối ưu hiệu năng, fix bug và đưa lên GitHub Pages.
\end{itemize}

\newpage
\section{Phân tích yêu cầu và kiến trúc hệ thống}
\subsection{Yêu cầu chức năng}
\begin{enumerate}[leftmargin=*]
    \item Nhận diện tư thế thời gian thực (Ngồi đúng, Cúi đầu, Vẹo lưng, Mắt quá gần).
    \item Cảnh báo tự động bằng âm thanh/popup sau 30s sai tư thế.
    \item Tích hợp đồng hồ Pomodoro timer.
    \item Lưu trữ phiên học, hệ thống Chuỗi lửa (Streak), Nhiệm vụ hàng ngày.
    \item Chat nhóm realtime, gửi file (hình ảnh, tệp tin).
\end{enumerate}

\subsection{Kiến trúc tổng quát}
Ứng dụng thiết kế theo kiến trúc Client-side (SPA tĩnh) kết hợp Backend-as-a-Service (Firebase).
\begin{figure}[H]
    \centering
    \fbox{\parbox[c][4cm][c]{0.8\textwidth}{\centering \textbf{Sơ đồ luồng xử lý:}\\ Webcam $\rightarrow$ TensorFlow.js (MoveNet) $\rightarrow$ Logic Phân loại $\rightarrow$ Giao diện (Cảnh báo/Thống kê) $\leftrightarrow$ Firebase (Lưu trữ/Chat)}}
    \caption{Kiến trúc xử lý phía Client và đồng bộ Firebase.}
\end{figure}

\section{Thiết kế chi tiết và thuật toán}
\subsection{Thuật toán phân loại tư thế (Heuristics)}
\begin{itemize}
    \item \textbf{Cúi đầu:} Tính tỷ lệ khoảng cách trục Y từ điểm Mũi đến Tâm 2 vai. Cân chỉnh tránh báo động giả, nếu tỷ lệ quá thấp $\Rightarrow$ Cảnh báo.
    \item \textbf{Vẹo lưng:} Đo góc nghiêng của đường nối 2 vai. Nếu $|angle| > 8^\circ \Rightarrow$ Cảnh báo.
    \item \textbf{Mắt quá gần:} Đo khoảng cách 2 tai. Nếu chiếm $> 32\%$ chiều ngang khung hình $\Rightarrow$ Cảnh báo.
\end{itemize}

\subsection{Kỹ thuật Làm mịn (Smoothing)}
Để tránh hiện tượng nhấp nháy, ứng dụng lưu 5 khung hình gần nhất và áp dụng thuật toán Majority Vote (Bầu chọn số đông) để đưa ra kết luận tư thế cuối cùng.

\section{Triển khai tính năng Mạng xã hội}
Hệ thống sử dụng Firebase Realtime Database và Storage:
\begin{itemize}
    \item \textbf{Chat realtime:} Dữ liệu tin nhắn cập nhật tức thì. Hỗ trợ gửi text, emoji, file base64 (file nhỏ) hoặc Storage (file lớn).
    \item \textbf{Phòng học ảo:} Theo dõi trạng thái Online/Offline của bạn bè.
\end{itemize}

\section{Kiểm thử và Đánh giá}
\subsection{Kết quả đo lường hiệu năng}
\begin{table}[H]
    \centering
    \begin{tabular}{|l|c|l|}
    \hline
    \textbf{Môi trường} & \textbf{FPS trung bình} & \textbf{Ghi chú} \\
    \hline
    Chrome (Laptop Core i5) & 28 -- 32 & Tối ưu tốt với MoveNet Lightning \\
    Firefox (Laptop) & 22 -- 28 & Thấp hơn một chút do WebGL renderer \\
    Thiết bị Di động (Mobile) & 15 -- 20 & Hoạt động mượt nếu đủ sáng \\
    \hline
    \end{tabular}
    \caption{Kết quả đo FPS thực tế}
\end{table}
Tỷ lệ nhận diện sai (False Positive) chỉ khoảng 8\% - 10\%, có thể khắc phục bằng cách tăng độ sáng phòng học.

\section{Tổng kết}
\subsection{Điểm mạnh}
\begin{itemize}
    \item Bảo mật 100\% hình ảnh (không upload camera lên mạng).
    \item Giao diện đẹp (V2), nhẹ, chạy trực tiếp trên web không cần cài đặt.
    \item Gamification (Trò chơi hóa) tạo động lực rất tốt.
\end{itemize}

\subsection{Hạn chế \& Hướng phát triển}
\begin{itemize}
    \item Hạn chế ở môi trường thiếu sáng hoặc khi có nhiều người lọt vào khung hình.
    \item Kế hoạch: Áp dụng các mô hình học sâu (Deep Learning) 3D Pose nếu cấu hình thiết bị người dùng trong tương lai đủ mạnh, bổ sung chế độ tự hiệu chỉnh (Calibration) cho từng người.
\end{itemize}

\newpage
\appendix
\section{Phụ lục: Code Snippet Vòng lặp Nhận diện}
\begin{lstlisting}[language=JavaScript]
async function predictionLoop(){
    if(!isRunning) return;
    const poses = await detector.estimatePoses(video);
    drawFrame(poses);
    if(poses.length > 0) {
        let rawPose = classifyPose(poses[0].keypoints);
        let finalPose = smoothPoseResult(rawPose);
        updateUI(finalPose);
    }
    requestAnimationFrame(predictionLoop);
}
\end{lstlisting}

\vspace{1cm}
\noindent \textbf{Người lập báo cáo:} Vương Lâm \\
\noindent \textbf{Ngày:} \today

\end{document}
